% SI-Methods Section S7q — Burnside's Lemma for Conformer Orbit Counting
%
% Proposed addition to si_methods.tex (located at
% /home/qmchem_max/draft_SMILES2xyz/si_methods.tex).
% NOT written to the draft repo per the worktree-isolation policy of
% this task brief.  Merge manually after review.
%
% Cross-references assumed already defined elsewhere in the SI:
%   \cref{si:polya}        — S7.x Pólya isomer counting
%   \cref{si:cp_pucker}    — S7.x Cremer-Pople ring conformers
%   \cref{si:grip_ensemble}— S7.x GRIP-ensemble overview

\subsection{S7q. Burnside's lemma for conformer orbit counting}
\label{si:burnside_conformer}

The Pólya enumeration on coordination-isomer topology
(\cref{si:polya}) and the Cremer-Pople canonical-state enumeration
on individual rings (\cref{si:cp_pucker}) together generate a
finite product space of candidate three-dimensional conformers,
\begin{equation}
  X \;=\; X_{\mathrm{iso}}
  \,\times\, \prod_{r=1}^{R} X_{\mathrm{ring}, r}
  \,\times\, X_{\mathrm{rot}},
  \label{eq:product_space}
\end{equation}
of cardinality
$\lvert X \rvert = n_{\mathrm{iso}} \cdot \prod_r m_r \cdot n_{\mathrm{rot}}$.
This raw product, however, in general over-counts because the
\emph{molecular point group} $G\subseteq O(3)$ of the assembled
complex acts non-trivially on $X$, identifying configurations that
differ only by a symmetry-equivalent relabelling of equivalent
ligands or arms.  For example, the $C_3$ axis of
$\Delta\text{-}[\mathrm{Co(en)}_3]^{3+}$ permutes the three ethylene-
diamine chelate rings, so the three distinct $\delta\delta\lambda$
assignments (depending on which ring carries the $\lambda$
puckering) represent a single physical conformer.

To recover the rigorous count of distinct conformers, we apply
\emph{Burnside's lemma} (also called the Cauchy--Frobenius
counting theorem):
\begin{equation}
  \lvert X / G \rvert
  \;=\;
  \frac{1}{\lvert G \rvert} \sum_{g \in G} \lvert \mathrm{Fix}(g) \rvert,
  \label{eq:burnside}
\end{equation}
where $\mathrm{Fix}(g) = \{ x \in X : g \cdot x = x\}$ is the set of
elements fixed by the group element $g$.  In our implementation
(\texttt{delfin/fffree/burnside\_conformer.py}), Eq.~\eqref{eq:burnside}
is realised in two complementary forms:
\begin{enumerate}
  \item For an \emph{explicit} candidate set $X$, we compute the
        induced permutation action of $G$ on the canonical-form
        index set --- each $g$ becomes a permutation of $\{1,\dots,
        \lvert X \rvert\}$ via inertial-frame alignment and
        atom-tuple lexicographic sort --- and apply
        Eq.~\eqref{eq:burnside} directly
        (\texttt{count\_distinct\_conformers}).
  \item For a \emph{formal} product space whose cardinalities are
        known but whose elements have not been materialised, we use
        the cyclic-action specialisation: when $C_k \subseteq G$
        permutes $kt$ equivalent ring slots, each filled by one of
        $m$ states, Pólya's cycle index gives
        \begin{equation}
          \frac{1}{k}\sum_{j=0}^{k-1} m^{\,t\,\gcd(j,k)},
          \label{eq:cycle_index}
        \end{equation}
        and multiplied with the unaffected factors yields the orbit
        count (\texttt{count\_orbits\_product\_space}).
\end{enumerate}

\paragraph{Worked example: $[\mathrm{Co(en)}_3]^{3+}$.}
With $n_{\mathrm{iso}} = 2$ ($\Delta, \Lambda$),
three en chelate rings each contributing two CP states
($\delta, \lambda$), no additional rotatable bonds
($n_{\mathrm{rot}} = 1$), and the $C_3$ axis permuting the three
rings, the raw product is $\lvert X \rvert = 2 \cdot 2^3 \cdot 1 = 16$.
Equation~\eqref{eq:cycle_index} gives
$\tfrac{1}{3}\!\left(2^3 + 2 + 2\right) = 4$ orbits on the ring
slots, and the full count
$\lvert X / C_3 \rvert = 2 \cdot 4 \cdot 1 = 8$ matches the textbook
enumeration $\{\Delta, \Lambda\} \times \{\delta\delta\delta,
\delta\delta\lambda, \delta\lambda\lambda, \lambda\lambda\lambda\}$.

\paragraph{Connection to existing layers.}
The three counting layers form a subgroup tower
\begin{equation*}
  G_{\mathrm{complex}} \;\supseteq\; G_{\mathrm{polyhedron}}
  \;\supseteq\; G_{\mathrm{ring},1} \times \cdots \times G_{\mathrm{torsion}},
\end{equation*}
and Eq.~\eqref{eq:burnside} restricted to each subgroup reproduces
the corresponding existing count: the Pólya isomer count of
\cref{si:polya} when restricted to $G_{\mathrm{polyhedron}}$
acting on the iso-factor, the Cremer-Pople ring conformer count of
\cref{si:cp_pucker} when restricted to $G_{\mathrm{ring},r}$.
The new \emph{conformer-level} layer adds the additional elements of
$G_{\mathrm{complex}}$ that interchange equivalent ligands, which is
the contribution that no competing tool currently provides
(ETKDG\cite{Riniker2015}, CREST\cite{Pracht2020},
OMEGA\cite{Hawkins2010}, molSimplify\cite{Ioannidis2016},
Frog2\cite{Miteva2010}, CSD Conformer Generator\cite{Cole2018}).

\paragraph{Determinism contract.}
The orbit count is independent of insertion order and floating-point
representation.  We enforce \texttt{PYTHONHASHSEED=0}, iterate over
explicit (not dict-derived) sequences, sort atom tuples
lexicographically after rounding coordinates to a tolerance
$\epsilon_{\mathrm{pos}} = 10^{-4}\,\text{\AA}$, and compare canonical
forms within an RMS tolerance $\epsilon_{\mathrm{rms}} = 0.05\,
\text{\AA}$.  The closed-form Eq.~\eqref{eq:cycle_index} is integer
arithmetic.  Two independent runs on the same input return
byte-identical orbit counts (test
\texttt{test\_determinism\_two\_runs\_identical}).

\paragraph{Completeness verdict.}
For reporting in
\cref{si:grip_ensemble}, we expose the function
\texttt{completeness\_report} which returns a dictionary with the
raw product, the Burnside orbit count, the emitted count, the
reduction factor, and a label
$\in \{\textsf{complete}, \textsf{mostly\_complete},
\textsf{incomplete}, \textsf{redundant}, \textsf{empty}\}$.  When
the emitted count agrees with the Burnside count to within $\pm 1$
(allowing for tolerance ties), the ensemble is \emph{provably
mathematically complete} on the chosen product space under the
detected point group.  This is the technical content of our
completeness claim.
